% =============================================================
%  RESUMEN TEÓRICO — Taller de Datos VI (1S-2026)
%  Cubre: Intro, Camino ML, Modelos Lineales, Árboles,
%         Ensambles I–III, Reducción de la Dimensionalidad
% =============================================================
\documentclass[12pt, a4paper]{article}

% --- Paquetes ---
\usepackage[utf8]{inputenc}
\usepackage[T1]{fontenc}
\usepackage[spanish]{babel}
\usepackage{amsmath, amssymb, amsthm}
\usepackage{geometry}
\usepackage{booktabs}
\usepackage{array}
\usepackage{graphicx}
\usepackage{xcolor}
\usepackage{hyperref}
\usepackage{enumitem}
\usepackage{tikz}
\usepackage{pgfplots}
\usepackage{algorithm}
\usepackage{algpseudocode}
\usepackage{mathtools}
\usepackage{bm}
\usepackage{tcolorbox}
\usepackage{multicol}
\usepackage{float}

\usetikzlibrary{shapes.geometric, arrows.meta, positioning, fit, calc, decorations.pathreplacing, matrix, backgrounds, babel}
\pgfplotsset{compat=1.18}

\geometry{top=2.5cm, bottom=2.5cm, left=2.5cm, right=2.5cm}

% --- Colores ---
\definecolor{azul}{RGB}{31, 78, 121}
\definecolor{verde}{RGB}{0, 112, 60}
\definecolor{rojo}{RGB}{180, 35, 24}
\definecolor{gris}{RGB}{230, 230, 230}
\definecolor{naranja}{RGB}{200, 100, 0}
\definecolor{morado}{RGB}{100, 20, 140}

% --- Cajas de definición ---
\tcbuselibrary{skins, breakable}
\newtcolorbox{defbox}[1]{
  colback=azul!8, colframe=azul!70, fonttitle=\bfseries,
  title={#1}, breakable, arc=3pt
}
\newtcolorbox{algbox}[1]{
  colback=verde!8, colframe=verde!70, fonttitle=\bfseries,
  title={#1}, breakable, arc=3pt
}
\newtcolorbox{notebox}{
  colback=naranja!10, colframe=naranja!70, breakable, arc=3pt
}

% --- Operadores matemáticos ---
\DeclareMathOperator{\sign}{sign}
\DeclareMathOperator{\Var}{Var}
\DeclareMathOperator{\Cov}{Cov}
\newcommand{\E}{\mathbb{E}}
\newcommand{\R}{\mathbb{R}}

% =============================================================
\begin{document}

\begin{titlepage}
	\centering
	\vspace*{2cm}
	{\Huge\bfseries\color{azul} Resumen \\[0.4em]
		Tecnología Digital VI\par}
	\vspace{0.5cm}
	{\large 1er Semestre 2026\par}
	\vspace{3cm}
	{\large\itshape Materias cubiertas: Introducción, Camino ML, Modelos Lineales,
		Árboles de Decisión, Ensambles I–III, Reducción de la Dimensionalidad\par}
\end{titlepage}

\tableofcontents
\newpage

% =============================================================
\section{Preprocesamiento y Validación}
% =============================================================

\subsection{Calidad de Datos}

\subsubsection{Valores Faltantes}
Existen tres mecanismos de ausencia:
\begin{itemize}
	\item \textbf{MCAR} (Missing Completely at Random): la ausencia no depende de ningún valor.
	\item \textbf{MAR} (Missing at Random): la ausencia depende de otras variables observadas.
	\item \textbf{MNAR} (Missing Not at Random): la ausencia depende del propio valor faltante.
\end{itemize}

\subsubsection{Outliers}
\begin{itemize}
	\item \textbf{Z-score}: $z_i = \dfrac{x_i - \mu}{\sigma}$; se considera outlier si $|z_i| > 3$.
	\item \textbf{Método IQR}: outlier si $x < Q_1 - 1.5\,\text{IQR}$ o $x > Q_3 + 1.5\,\text{IQR}$.
\end{itemize}

\subsubsection{Desbalance de Clases}
\begin{itemize}
	\item \textbf{Oversampling}: duplicar muestras de la clase minoritaria.
	\item \textbf{Undersampling}: eliminar muestras de la clase mayoritaria.
	\item \textbf{SMOTE}: generar muestras sintéticas interpolando vecinos de la clase minoritaria.
	\item \textbf{Class weights}: penalizar más la clasificación errónea de la clase minoritaria.
\end{itemize}

\subsubsection{Data Leakage}
\begin{notebox}
	El \textbf{data leakage} ocurre cuando información del conjunto de test (o del futuro)
	se filtra al entrenamiento. Consecuencia: métricas optimistas que no se reproducen en
	producción. Prevención: ajustar transformadores \emph{solo} con datos de entrenamiento.
\end{notebox}

\subsection{Escalado de Features}
\[
	\text{MinMax: } x' = \frac{x - x_{\min}}{x_{\max} - x_{\min}}, \qquad
	\text{Estándar: } x' = \frac{x - \mu}{\sigma}, \qquad
	\text{Robusto: } x' = \frac{x - \tilde{x}}{\text{IQR}}
\]

\subsection{Selección de Features}
\begin{center}
	\begin{tabular}{lll}
		\toprule
		\textbf{Tipo} & \textbf{Método}                          & \textbf{Característica}        \\
		\midrule
		Filter        & Correlación, $\chi^2$, información mutua & Independiente del modelo       \\
		Wrapper       & Forward/Backward selection               & Evalúa subconjuntos con modelo \\
		Embedded      & L1 (Lasso), importancia en árboles       & Integrado al entrenamiento     \\
		\bottomrule
	\end{tabular}
\end{center}

\subsection{K-Fold Cross-Validation}

\begin{figure}[H]
	\centering
	\begin{tikzpicture}[scale=0.8]
		\foreach \fold in {1,...,5} {
				\foreach \part in {1,...,5} {
						\ifnum\fold=\part
							\fill[rojo!60] ({(\part-1)*2}, {-((\fold-1)*0.7)}) rectangle ({(\part)*2}, {-((\fold-1)*0.7)+0.55});
						\else
							\fill[azul!30] ({(\part-1)*2}, {-((\fold-1)*0.7)}) rectangle ({(\part)*2}, {-((\fold-1)*0.7)+0.55});
						\fi
						\draw[gray!50] ({(\part-1)*2}, {-((\fold-1)*0.7)}) rectangle ({(\part)*2}, {-((\fold-1)*0.7)+0.55});
					}
				\node[right, font=\small] at (10.2, {-((\fold-1)*0.7)+0.28}) {Fold \fold};
			}
		\node[azul!70, font=\small\bfseries] at (5, 0.9) {Entrenamiento};
		\node[rojo!70, font=\small\bfseries] at (5, -3.6) {Validación (rojo)};
	\end{tikzpicture}
	\caption{5-Fold Cross-Validation: en cada iteración un fold actúa como validación.}
\end{figure}

\[
	\hat{R}_{CV} = \frac{1}{k}\sum_{j=1}^{k} \text{score}_j
\]

% =============================================================
\section{Métricas de Evaluación}
% =============================================================

\subsection{Regresión}

\begin{defbox}{Métricas de Regresión}
	\begin{align}
		\text{MSE}   & = \frac{1}{n}\sum_{i=1}^{n}(y_i - \hat{y}_i)^2                                \\[4pt]
		\text{RMSE}  & = \sqrt{\text{MSE}}                                                           \\[4pt]
		\text{MAE}   & = \frac{1}{n}\sum_{i=1}^{n}|y_i - \hat{y}_i|                                  \\[4pt]
		\text{MAPE}  & = \frac{100}{n}\sum_{i=1}^{n}\frac{|y_i - \hat{y}_i|}{y_i}                    \\[4pt]
		\text{RMSLE} & = \sqrt{\frac{1}{n}\sum_{i=1}^{n}\bigl(\log(y_i+1)-\log(\hat{y}_i+1)\bigr)^2} \\[4pt]
		R^2          & = 1 - \frac{\sum(y_i-\hat{y}_i)^2}{\sum(y_i-\bar{y})^2}
	\end{align}
\end{defbox}

\subsection{Clasificación}

\subsubsection{Matriz de Confusión}

\begin{center}
	\begin{tabular}{lcc}
		\toprule
		                       & \textbf{Pred. Positivo} & \textbf{Pred. Negativo} \\
		\midrule
		\textbf{Real Positivo} & \textcolor{verde}{TP}   & \textcolor{rojo}{FN}    \\
		\textbf{Real Negativo} & \textcolor{rojo}{FP}    & \textcolor{verde}{TN}   \\
		\bottomrule
	\end{tabular}
\end{center}

\begin{defbox}{Métricas de Clasificación}
	\[
		\text{Accuracy} = \frac{TP+TN}{TP+TN+FP+FN}, \quad
		\text{Precision} = \frac{TP}{TP+FP}, \quad
		\text{Recall} = \frac{TP}{TP+FN}
	\]
	\[
		F_1 = \frac{2 \cdot \text{Precision} \cdot \text{Recall}}{\text{Precision}+\text{Recall}}, \quad
		\text{TPR} = \frac{TP}{TP+FN}, \quad
		\text{FPR} = \frac{FP}{FP+TN}
	\]
\end{defbox}

\subsubsection{Curva ROC y AUC}
La curva ROC grafica TPR vs.\ FPR para distintos umbrales de clasificación.

\begin{figure}[H]
	\centering
	\begin{tikzpicture}
		\begin{axis}[
				width=7cm, height=6cm,
				xlabel={FPR}, ylabel={TPR},
				xmin=0, xmax=1, ymin=0, ymax=1,
				xtick={0,0.2,...,1}, ytick={0,0.2,...,1},
				grid=both, grid style={gray!20},
				title={Curva ROC}
			]
			\addplot[azul, thick, smooth] coordinates {
					(0,0)(0.05,0.45)(0.1,0.6)(0.2,0.75)(0.3,0.85)(0.5,0.93)(0.7,0.97)(1,1)
				};
			\addplot[dashed, gray] coordinates {(0,0)(1,1)};
			\addlegendentry{Modelo}
			\addlegendentry{Azar}
		\end{axis}
	\end{tikzpicture}
	\caption{AUC-ROC: área bajo la curva. AUC = 0.5 equivale a azar; AUC = 1 es perfecto.}
\end{figure}

\subsubsection{Curva Precision-Recall}
Más informativa que ROC cuando las clases están muy desbalanceadas; se focaliza en la
clase positiva.

% =============================================================
\section{Modelos Lineales}
% =============================================================

\subsection{Regresión Lineal}

\begin{defbox}{Modelo}
	\[
		y_i = \beta_0 + \beta_1 x_{i1} + \cdots + \beta_p x_{ip} + \varepsilon_i
		\quad\Longleftrightarrow\quad \mathbf{y} = \mathbf{X}\bm{\beta} + \bm{\varepsilon}
	\]
	donde $\varepsilon_i \overset{iid}{\sim} \mathcal{N}(0,\sigma^2)$.
\end{defbox}

\subsubsection{Función de Pérdida y Solución Analítica}
\[
	J(\bm{\beta}) = \|\mathbf{y} - \mathbf{X}\bm{\beta}\|^2
	= \sum_{i=1}^{n}(y_i - \hat{y}_i)^2
\]
\textbf{Ecuaciones normales}:
\[
	\boxed{\hat{\bm{\beta}} = (\mathbf{X}^\top \mathbf{X})^{-1}\mathbf{X}^\top\mathbf{y}}
\]

\subsection{Regresión Logística}

\subsubsection{De probabilidades a log-odds}
\[
	\text{odds} = \frac{p}{1-p}, \qquad
	\text{logit}(p) = \ln\!\left(\frac{p}{1-p}\right)
\]

\subsubsection{Función Sigmoide}
\[
	\sigma(z) = \frac{1}{1+e^{-z}} \in (0,1)
\]

\begin{figure}[H]
	\centering
	\begin{tikzpicture}
		\begin{axis}[
				width=8cm, height=5cm,
				xlabel={$z$}, ylabel={$\sigma(z)$},
				xmin=-6, xmax=6, ymin=-0.05, ymax=1.05,
				ytick={0,0.5,1}, grid=both, grid style={gray!20},
				title={Función Sigmoide}
			]
			\addplot[azul, thick, domain=-6:6, samples=100] {1/(1+exp(-x))};
			\addplot[dashed, gray] coordinates {(-6,0.5)(6,0.5)};
		\end{axis}
	\end{tikzpicture}
	\caption{La sigmoide mapea cualquier valor real a $(0,1)$, interpretable como probabilidad.}
\end{figure}

\subsubsection{Modelo y Pérdida}
\[
	P(Y=1\mid\mathbf{x}) = \sigma(\bm{\beta}^\top\mathbf{x})
\]
\textbf{Log-loss (cross-entropía binaria)}:
\[
	J(\bm{\beta}) = -\frac{1}{n}\sum_{i=1}^{n}
	\bigl[y_i\log\hat{y}_i + (1-y_i)\log(1-\hat{y}_i)\bigr]
\]

\subsection{Regularización}

La regularización agrega un término de penalización a la función de pérdida para
controlar la complejidad del modelo y evitar el sobreajuste.

\begin{defbox}{Formas de Regularización}
	\begin{align}
		\text{Lasso (L1):} \quad  & J_{\text{Lasso}}(\mathbf{w}) = J(\mathbf{w}) + \lambda\sum_j |w_j| \\[4pt]
		\text{Ridge (L2):} \quad  & J_{\text{Ridge}}(\mathbf{w}) = J(\mathbf{w}) + \lambda\sum_j w_j^2 \\[4pt]
		\text{Elastic Net:} \quad & J_{\text{EN}}(\mathbf{w}) = J(\mathbf{w})
		+ \lambda_1\sum_j|w_j| + \lambda_2\sum_j w_j^2
	\end{align}
\end{defbox}

\begin{figure}[H]
	\centering
	\begin{tikzpicture}[scale=1.1]
		% L1 diamond
		\begin{scope}[xshift=-3.5cm]
			\draw[fill=azul!15, draw=azul, thick] (0,1)--(1,0)--(0,-1)--(-1,0)--cycle;
			\draw[azul!40, thick] (0.6,0.6) ellipse (1.2cm and 0.8cm);
			\fill[rojo] (-1,0) circle (2pt);
			\node[below, font=\small] at (0,-1.4) {L1: solución esparsa};
		\end{scope}
		% L2 circle
		\begin{scope}[xshift=3.5cm]
			\draw[fill=verde!15, draw=verde, thick] (0,0) circle (1cm);
			\draw[verde!40, thick] (0.7,0.7) ellipse (1.2cm and 0.8cm);
			\fill[rojo] ({cos(45)},{sin(45)}) circle (2pt);
			\node[below, font=\small] at (0,-1.4) {L2: coeficientes pequeños};
		\end{scope}
		\node at (0,0) {\large vs.};
	\end{tikzpicture}
	\caption{L1 tiende a producir coeficientes exactamente cero (selección de features);
		L2 encoge todos los coeficientes pero raramente a cero.}
\end{figure}

\begin{center}
	\begin{tabular}{lll}
		\toprule
		\textbf{Propiedad}    & \textbf{L1 (Lasso)}        & \textbf{L2 (Ridge)} \\
		\midrule
		Coeficientes          & Exactamente 0 (esparsidad) & Pequeños, $\neq 0$  \\
		Selección de features & Sí                         & No                  \\
		Multicolinealidad     & Selecciona una             & Distribuye el peso  \\
		Solución              & No analítica               & Analítica           \\
		\bottomrule
	\end{tabular}
\end{center}

% =============================================================
\section{Árboles de Decisión}
% =============================================================

\subsection{Intuición Geométrica}

Un árbol de decisión particiona el espacio de features en regiones rectangulares
(hiperplanos perpendiculares a los ejes). Dentro de cada hoja, la predicción es constante.

\begin{figure}[H]
	\centering
	\begin{tikzpicture}[scale=0.9]
		% Partición 2D
		\begin{scope}[xshift=-5cm]
			\draw[->] (0,0) -- (4.2,0) node[right]{\small $x_1$};
			\draw[->] (0,0) -- (0,4.2) node[above]{\small $x_2$};
			\draw[azul, thick] (2.5,0) -- (2.5,4);
			\draw[verde, thick] (0,2) -- (2.5,2);
			\draw[rojo, thick] (0,3.2) -- (2.5,3.2);
			\node at (1.2,1)  {\small $R_1$};
			\node at (1.2,2.6){\small $R_2$};
			\node at (1.2,3.6){\small $R_3$};
			\node at (3.3,2)  {\small $R_4$};
			\node[below, font=\footnotesize] at (2,-.4) {Partición del espacio};
		\end{scope}
		% Árbol correspondiente
		\begin{scope}[xshift=1cm]
			\tikzset{nodo/.style={rectangle,rounded corners,draw=azul,fill=azul!10,
						minimum width=2cm, minimum height=0.7cm, font=\footnotesize}}
			\tikzset{hoja/.style={rectangle,draw=verde,fill=verde!10,
						minimum width=1.4cm, minimum height=0.7cm, font=\footnotesize}}
			\node[nodo] (r) at (1.5,3.5) {$x_1 \leq 2.5$?};
			\node[nodo] (l) at (0,2)     {$x_2 \leq 2$?};
			\node[hoja] (r4) at (3,2)    {$R_4$};
			\node[nodo] (ll) at (-0.8,0.5){$x_2\leq 3.2$?};
			\node[hoja] (r1) at (0.8,0.5) {$R_1$};
			\node[hoja] (r2) at (-1.6,-1) {$R_2$};
			\node[hoja] (r3) at (0,-1)    {$R_3$};
			\draw[-Stealth] (r) -- node[left,font=\tiny]{sí} (l);
			\draw[-Stealth] (r) -- node[right,font=\tiny]{no} (r4);
			\draw[-Stealth] (l) -- node[left,font=\tiny]{sí} (ll);
			\draw[-Stealth] (l) -- node[right,font=\tiny]{no} (r1);
			\draw[-Stealth] (ll) -- node[left,font=\tiny]{sí}  (r2);
			\draw[-Stealth] (ll) -- node[right,font=\tiny]{no} (r3);
		\end{scope}
	\end{tikzpicture}
	\caption{Un árbol divide recursivamente el espacio de features; cada hoja corresponde
		a una región rectangular con predicción constante.}
\end{figure}

\subsection{Funciones de Impureza}

Para \textbf{regresión}, se minimiza el MSE ponderado del nodo:
\[
	\text{MSE}_m = \frac{1}{N_m}\sum_{i \in m}(y_i - \bar{y}_m)^2
\]

Para \textbf{clasificación} se emplean:
\begin{defbox}{Gini e Entropía}
	\[
		\text{Gini}_m = 1 - \sum_{k=1}^{K} p_{mk}^2, \qquad
		H_m = -\sum_{k=1}^{K} p_{mk}\log_2(p_{mk})
	\]
	donde $p_{mk}$ es la proporción de clase $k$ en el nodo $m$.
	Ambas valen 0 para nodos puros y alcanzan su máximo para distribuciones uniformes.
\end{defbox}

\subsection{Criterio de División}

Se elige el feature $j$ y umbral $t$ que minimizan la impureza ponderada de los hijos:
\[
	J(j, t) = \frac{N_{\text{izq}}}{N_m}\,I_{\text{izq}} + \frac{N_{\text{der}}}{N_m}\,I_{\text{der}}
\]

\begin{algbox}{Algoritmo CART (greedy)}
	\begin{algorithmic}[1]
		\For{cada feature $j$ y cada umbral $t$}
		\State Calcular $J(j,t)$
		\EndFor
		\State Dividir con $(j^*, t^*) = \arg\min_{j,t} J(j,t)$
		\State Repetir recursivamente en cada hijo hasta criterio de parada
	\end{algorithmic}
\end{algbox}

\subsection{Regularización del Árbol}

\begin{center}
	\begin{tabular}{lll}
		\toprule
		\textbf{Hiperparámetro}      & \textbf{Efecto}                 & \textbf{Rango típico}  \\
		\midrule
		\texttt{max\_depth}          & Profundidad máxima              & 3–10                   \\
		\texttt{min\_samples\_split} & Mín. muestras para dividir      & 2–20                   \\
		\texttt{min\_samples\_leaf}  & Mín. muestras por hoja          & 1–10                   \\
		\texttt{max\_features}       & Features consideradas por split & $\sqrt{p}$, $\log_2 p$ \\
		\bottomrule
	\end{tabular}
\end{center}

\subsection{Importancia de Features}

\[
	\text{Imp}(j) \propto \sum_{\text{splits en }j}
	\frac{N_m}{N}\bigl(I_m - \tfrac{N_{\text{izq}}}{N_m}I_{\text{izq}}
	- \tfrac{N_{\text{der}}}{N_m}I_{\text{der}}\bigr)
\]
Se pondera por la proporción de muestras que pasan por cada split y se normaliza para que
sumen 1.

\begin{notebox}
	\textbf{Sesgo de importancia}: los árboles tienden a sobreestimar la importancia de
	features con muchas categorías o alta cardinalidad. Para corregir esto se puede usar
	importancia por permutación.
\end{notebox}

% =============================================================
\section{Ensambles I: Voting y Bagging}
% =============================================================

\subsection{¿Qué es un Ensamble?}

\begin{defbox}{Ensamble}
	Combinar múltiples \textbf{modelos base} (learners) para producir una predicción
	más robusta que cualquier modelo individual.
\end{defbox}

\subsection{Votación (Voting)}

\paragraph{Clasificación:} predicción = moda de los votos.
\paragraph{Regresión:} predicción = promedio de las predicciones.

\textbf{¿Por qué funciona?}
Si $n$ clasificadores son independientes con exactitud $p > 0.5$ cada uno, la probabilidad
de que la mayoría se equivoque decrece exponencialmente con $n$:
\[
	P(\text{error del ensamble}) = \sum_{k=\lceil n/2\rceil}^{n}\binom{n}{k}(1-p)^k p^{n-k}
\]

\subsection{Varianza, Sesgo y Diversidad}

Para modelos \emph{no correlacionados} con varianza $\sigma^2$:
\[
	\Var\!\left(\frac{1}{B}\sum_{b=1}^{B}\hat{f}_b\right)
	= \frac{\sigma^2}{B} \xrightarrow{B\to\infty} 0
\]
En la práctica los modelos están correlacionados ($\rho$):
\[
	\Var(\bar{f}) = \rho\,\sigma^2 + \frac{1-\rho}{B}\,\sigma^2
\]
La clave es \textbf{reducir $\rho$}: modelos más diversos bajan más la varianza.

\subsection{Bootstrap y Bagging}

\begin{defbox}{Bootstrap}
	Dado un conjunto de $n$ muestras, seleccionar \emph{con reemplazo} $n$ muestras
	aleatoriamente. En promedio, el $\approx 63.2\%$ de las muestras originales aparece
	al menos una vez; el $\approx 36.8\%$ queda fuera (\emph{Out-of-Bag}, OOB).
\end{defbox}

\begin{algbox}{Bagging (Bootstrap Aggregating)}
	\begin{algorithmic}[1]
		\For{$b = 1$ to $B$}
		\State Generar muestra bootstrap $\mathcal{D}_b$ a partir de $\mathcal{D}$
		\State Entrenar modelo $\hat{f}_b$ en $\mathcal{D}_b$
		\EndFor
		\State \Return $\hat{f}(x) = \frac{1}{B}\sum_b\hat{f}_b(x)$ (o voto mayoritario)
	\end{algorithmic}
\end{algbox}

\paragraph{Validación OOB:} cada muestra es OOB para $\approx 37\%$ de los modelos;
se puede calcular un error de validación sin usar un conjunto separado.

\begin{figure}[H]
	\centering
	\begin{tikzpicture}[scale=0.75]
		\tikzset{ds/.style={rectangle, rounded corners, draw=azul, fill=azul!10,
					text width=2cm, align=center, minimum height=0.9cm, font=\footnotesize}}
		\tikzset{model/.style={rectangle, rounded corners, draw=verde, fill=verde!10,
					text width=1.6cm, align=center, minimum height=0.9cm, font=\footnotesize}}
		\node[ds] (D) at (0,3) {Dataset $\mathcal{D}$};
		\foreach \b/\x in {1/-4, 2/-1.3, 3/1.4, 4/4.1} {
				\node[ds] (b\b) at (\x, 1.5) {Bootstrap $\mathcal{D}_\b$};
				\node[model] (m\b) at (\x, 0) {Modelo $\hat{f}_\b$};
				\draw[-Stealth, azul!50] (D) -- (b\b);
				\draw[-Stealth, verde!50] (b\b) -- (m\b);
			}
		\node[rectangle, rounded corners, draw=rojo, fill=rojo!10,
			text width=3cm, align=center, minimum height=0.9cm, font=\footnotesize\bfseries]
		(agg) at (0,-1.5) {Agregación};
		\foreach \b in {1,...,4} \draw[-Stealth, rojo!50] (m\b) -- (agg);
	\end{tikzpicture}
	\caption{Esquema de Bagging: $B$ modelos entrenados en paralelo sobre muestras bootstrap.}
\end{figure}

% =============================================================
\section{Ensambles II: Random Forest e Interpretación}
% =============================================================

\subsection{Random Forest}

\begin{defbox}{Random Forest}
	Bagging aplicado a árboles de decisión con una fuente adicional de aleatoriedad:
	en cada split solo se evalúan $m$ features seleccionados al azar
	(típicamente $m = \sqrt{p}$ para clasificación, $m = p/3$ para regresión).
\end{defbox}

\paragraph{Efecto:} reduce la correlación entre los árboles $\Rightarrow$ mayor reducción
de varianza del ensamble.

\begin{figure}[H]
	\centering
	\begin{tikzpicture}
		\begin{axis}[
				width=9cm, height=5.5cm,
				xlabel={Número de árboles $B$},
				ylabel={Error OOB},
				xmin=1, xmax=200, ymin=0, ymax=0.35,
				grid=both, grid style={gray!15},
				title={Error OOB vs Número de Árboles}
			]
			\addplot[azul, thick, smooth] table[row sep=\\] {
					x y\\
					1 0.32\\ 5 0.22\\ 10 0.17\\ 20 0.14\\ 30 0.125\\
					50 0.11\\ 80 0.102\\ 120 0.098\\ 180 0.095\\ 200 0.094\\
				};
			\addplot[rojo, dashed, thick, smooth] table[row sep=\\] {
					x y\\
					1 0.28\\ 5 0.18\\ 10 0.14\\ 20 0.11\\ 30 0.1\\
					50 0.09\\ 80 0.085\\ 120 0.082\\ 180 0.08\\ 200 0.079\\
				};
			\addlegendentry{Árbol simple (bagging)}
			\addlegendentry{Random Forest}
		\end{axis}
	\end{tikzpicture}
	\caption{El error estabiliza después de cierto número de árboles (retornos decrecientes).}
\end{figure}

\subsection{Dependencia Parcial (PDP)}

La \textbf{función de dependencia parcial} mide el efecto marginal promedio del feature
$x_s$ sobre la predicción, promediando sobre todos los demás features:
\[
	\hat{f}_s(x_s) = \frac{1}{n}\sum_{i=1}^{n}\hat{f}(x_s,\, \mathbf{x}_{-s}^{(i)})
\]

\subsection{Individual Conditional Expectation (ICE)}

Similar al PDP pero se grafica una curva por muestra (sin promediar).
La media de todas las curvas ICE coincide con el PDP.

\begin{figure}[H]
	\centering
	\begin{tikzpicture}
		\begin{axis}[
				width=9cm, height=5cm,
				xlabel={Valor de feature $x_s$},
				ylabel={Predicción},
				grid=both, grid style={gray!15},
				title={ICE (gris) y PDP (azul grueso)}
			]
			% Curvas ICE (simuladas)
			\foreach \offset in {-0.3,-0.15,0,0.15,0.3}{
					\addplot[gray!50, thin] expression[domain=0:5, samples=30]
						{0.5*x + \offset*sin(deg(x*3)) + \offset + 1};
				}
			% PDP
			\addplot[azul, very thick] expression[domain=0:5, samples=50] {0.5*x + 1};
		\end{axis}
	\end{tikzpicture}
	\caption{Las curvas ICE revelan heterogeneidad en el efecto del feature.}
\end{figure}

% =============================================================
\section{Ensambles III: Boosting y Gradient Boosting}
% =============================================================

\subsection{Filosofía del Boosting}

\begin{defbox}{Boosting}
	Construir un modelo fuerte a partir de \textbf{learners débiles} entrenados
	\emph{secuencialmente}: cada modelo nuevo se concentra en los errores cometidos
	por el ensamble anterior.
\end{defbox}

\begin{figure}[H]
	\centering
	\begin{tikzpicture}[scale=0.8]
		\tikzset{m/.style={circle, draw=azul, fill=azul!15, minimum size=1cm, font=\small}}
		\tikzset{plus/.style={circle, draw=verde, fill=verde!15, minimum size=0.5cm, font=\small\bfseries}}
		\node[m] (m1) at (0,0) {$h_1$};
		\node[m] (m2) at (2.5,0) {$h_2$};
		\node[m] (m3) at (5,0) {$h_3$};
		\node at (7,0) {$\cdots$};
		\node[m] (mM) at (9,0) {$h_M$};
		\node[rectangle, draw=rojo, fill=rojo!10, minimum width=2.5cm, font=\small,
			align=center] (F) at (4.5,-2) {$F(x) = \sum_m \alpha_m h_m(x)$};
		\foreach \m in {m1,m2,m3,mM} \draw[-Stealth,rojo!60] (\m) -- (F);
		\draw[-Stealth, dashed, naranja] (m1) -- node[above,font=\tiny]{residuos} (m2);
		\draw[-Stealth, dashed, naranja] (m2) -- node[above,font=\tiny]{residuos} (m3);
		\draw[-Stealth, dashed, naranja] (m3) -- node[above,font=\tiny]{} (7,0);
		\draw[-Stealth, dashed, naranja] (7,0) -- node[above,font=\tiny]{residuos} (mM);
	\end{tikzpicture}
	\caption{Esquema de Boosting: cada learner débil aprende de los errores del anterior.}
\end{figure}

\subsection{AdaBoost}

\begin{algbox}{AdaBoost (clasificación binaria)}
	\begin{algorithmic}[1]
		\State Inicializar pesos: $w_i = 1/n$, $\forall i$
		\For{$m = 1$ to $M$}
		\State Entrenar clasificador $h_m$ con pesos $\{w_i\}$
		\State Calcular error ponderado: $\varepsilon_m = \sum_{h_m(x_i)\neq y_i} w_i$
		\State Calcular peso del modelo: $\alpha_m = \frac{1}{2}\ln\!\dfrac{1-\varepsilon_m}{\varepsilon_m}$
		\State Actualizar pesos:
		\[
			w_i \leftarrow w_i \cdot \exp\!\bigl(-\alpha_m\, y_i\, h_m(x_i)\bigr)
		\]
		\State Renormalizar: $w_i \leftarrow w_i / \sum_j w_j$
		\EndFor
		\State \Return $H(x) = \sign\!\left(\sum_m \alpha_m h_m(x)\right)$
	\end{algorithmic}
\end{algbox}

\paragraph{Intuición:} si $\varepsilon_m < 0.5$ entonces $\alpha_m > 0$; muestras mal
clasificadas reciben mayor peso en la siguiente iteración.

\subsection{Gradient Boosting}

\subsubsection{Algoritmo General}

\begin{algbox}{Gradient Boosting (Friedman, 1999)}
	\begin{algorithmic}[1]
		\State Inicializar: $F_0(x) = \arg\min_c \sum_i L(y_i, c)$
		\For{$m = 1$ to $M$}
		\State Calcular pseudo-residuos (gradiente negativo):
		\[
			r_{im} = -\frac{\partial L(y_i, F(x_i))}{\partial F(x_i)}\bigg|_{F=F_{m-1}}
		\]
		\State Entrenar árbol de regresión $h_m$ sobre $\{(x_i, r_{im})\}$
		\State Determinar paso óptimo:
		$\gamma_m = \arg\min_\gamma \sum_i L\bigl(y_i, F_{m-1}(x_i) + \gamma\, h_m(x_i)\bigr)$
		\State Actualizar: $F_m(x) = F_{m-1}(x) + \eta\,\gamma_m\, h_m(x)$
		\EndFor
		\State \Return $F_M(x)$
	\end{algorithmic}
\end{algbox}

donde $\eta \in (0,1]$ es la \textbf{tasa de aprendizaje} (shrinkage).

\subsubsection{Conexión con Descenso de Gradiente}

El gradient boosting es equivalente al \emph{descenso de gradiente en el espacio funcional}:
cada árbol aproxima el gradiente negativo de la pérdida respecto de la predicción actual.

\begin{center}
	\begin{tabular}{ll}
		\toprule
		\textbf{Función de pérdida} $L(y, F)$ & \textbf{Pseudo-residuo} $r$ \\
		\midrule
		$\tfrac{1}{2}(y-F)^2$ (regresión)     & $y - F$                     \\
		$\log(1+e^{-2yF})$ (clasificación)    & $\dfrac{2y}{1+e^{2yF}}$     \\
		$|y-F|$ (MAE)                         & $\sign(y-F)$                \\
		\bottomrule
	\end{tabular}
\end{center}

\subsubsection{Parámetros Clave}

\begin{center}
	\begin{tabular}{lll}
		\toprule
		\textbf{Parámetro}         & \textbf{Efecto}                      & \textbf{Rango típico} \\
		\midrule
		$M$ (n\_estimators)        & Número de árboles                    & 100–3000              \\
		$\eta$ (learning\_rate)    & Tamaño del paso; menor = más robusto & 0.01–0.3              \\
		\texttt{max\_depth}        & Complejidad de cada árbol            & 3–8                   \\
		\texttt{subsample}         & Fracción de filas por árbol          & 0.5–1.0               \\
		\texttt{colsample\_bytree} & Fracción de columnas por árbol       & 0.5–1.0               \\
		\bottomrule
	\end{tabular}
\end{center}

\begin{notebox}
	\textbf{Trade-off}: $\eta$ pequeño requiere $M$ grande; combinados producen mayor
	regularización implícita y mejor generalización, a costa de tiempo de entrenamiento.
\end{notebox}

\subsection{XGBoost: Objetivo Regularizado}

XGBoost extiende gradient boosting con una función objetivo regularizada explícita:
\[
	\mathcal{L} = \underbrace{\sum_{i=1}^n L(y_i, \hat{y}_i)}_{\text{pérdida}} +
	\underbrace{\sum_{k=1}^K \Omega(f_k)}_{\text{regularización}}
\]
donde:
\[
	\Omega(f_k) = \gamma\, T_k + \frac{\lambda}{2}\sum_{j=1}^{T_k} w_j^2
\]
$T_k$ = número de hojas, $w_j$ = valor en la hoja $j$, $\gamma$ y $\lambda$ son
hiperparámetros de regularización.

\subsubsection{Expansión de Taylor de Segundo Orden}
Para optimizar eficientemente el objetivo:
\[
	\mathcal{L}^{(m)} \approx \sum_i \left[
		g_i f_m(x_i) + \frac{1}{2} h_i f_m(x_i)^2
		\right] + \Omega(f_m)
\]
donde $g_i = \partial_{\hat{y}} L(y_i, \hat{y})$ y $h_i = \partial^2_{\hat{y}} L(y_i, \hat{y})$.

El valor óptimo en cada hoja $j$ es:
\[
	w_j^* = -\frac{\sum_{i\in j} g_i}{\sum_{i\in j} h_i + \lambda}
\]

\subsection{LightGBM y CatBoost}

\begin{center}
	\begin{tabular}{llll}
		\toprule
		                     & \textbf{XGBoost}  & \textbf{LightGBM}     & \textbf{CatBoost} \\
		\midrule
		Crecimiento árbol    & Por nivel         & Por hoja              & Simétrico         \\
		Velocidad            & Moderada          & Alta                  & Alta              \\
		Memoria              & Moderada          & Baja                  & Moderada          \\
		Features categóricas & Requiere encoding & Nativo                & Nativo            \\
		Regularización       & Explícita         & Implícita + explícita & Implícita         \\
		\bottomrule
	\end{tabular}
\end{center}

\begin{figure}[H]
	\centering
	\begin{tikzpicture}[scale=0.9,
			node/.style={circle, draw=azul!60, fill=azul!10, minimum size=0.7cm, font=\tiny},
			leaf/.style={rectangle, draw=verde!60, fill=verde!10, minimum width=0.7cm,
					minimum height=0.5cm, font=\tiny}
		]
		% Crecimiento por nivel (XGBoost)
		\begin{scope}[xshift=-4.5cm]
			\node[node] (r) at (0,2.5) {};
			\node[node] (l1) at (-1,1.5) {}; \node[node] (r1) at (1,1.5) {};
			\node[leaf] (l2) at (-1.5,0.5) {}; \node[leaf] (r2) at (-0.5,0.5) {};
			\node[leaf] (l3) at (0.5,0.5)  {}; \node[leaf] (r3) at (1.5,0.5)  {};
			\draw[-Stealth,azul!40] (r)--(l1); \draw[-Stealth,azul!40] (r)--(r1);
			\draw[-Stealth,azul!40] (l1)--(l2); \draw[-Stealth,azul!40] (l1)--(r2);
			\draw[-Stealth,azul!40] (r1)--(l3); \draw[-Stealth,azul!40] (r1)--(r3);
			\node[font=\footnotesize, below] at (0,-0.1) {Por nivel};
		\end{scope}
		% Crecimiento por hoja (LightGBM)
		\begin{scope}[xshift=3cm]
			\node[node] (r) at (0,2.5) {};
			\node[node] (l1) at (-1,1.5) {}; \node[leaf] (r1) at (1,1.5) {};
			\node[node] (l2) at (-1.5,0.5) {}; \node[leaf] (r2) at (-0.5,0.5) {};
			\node[leaf] (l3) at (-2,-0.5) {}; \node[leaf] (r3) at (-1,-0.5) {};
			\draw[-Stealth,verde!60] (r)--(l1); \draw[-Stealth,verde!60] (r)--(r1);
			\draw[-Stealth,verde!60] (l1)--(l2); \draw[-Stealth,verde!60] (l1)--(r2);
			\draw[-Stealth,verde!60] (l2)--(l3); \draw[-Stealth,verde!60] (l2)--(r3);
			\node[font=\footnotesize, below] at (-0.5,-0.9) {Por hoja (leaf-wise)};
		\end{scope}
	\end{tikzpicture}
	\caption{XGBoost expande el árbol nivel por nivel; LightGBM expande la hoja con mayor
		reducción de pérdida (más eficiente pero potencialmente más overfit).}
\end{figure}

% =============================================================
\section{Reducción de la Dimensionalidad}
% =============================================================

\subsection{Motivación: La Maldición de la Dimensionalidad}

\begin{notebox}
	En espacios de alta dimensión el volumen crece exponencialmente: los puntos se vuelven
	\emph{equidistantes} y los algoritmos basados en distancia (kNN, SVM radial, etc.)
	pierden efectividad. Además, visualizar y comprender datos con miles de features es
	prácticamente imposible.
\end{notebox}

\subsection{Análisis de Componentes Principales (PCA)}

\subsubsection{Objetivo}
Encontrar una proyección lineal $\mathbf{z} = \mathbf{W}^\top \mathbf{x}$
($\mathbf{W} \in \R^{p\times k}$, $k < p$) que maximice la varianza retenida,
con columnas de $\mathbf{W}$ ortonormales.

\subsubsection{Herramientas Estadísticas Previas}

\paragraph{Varianza} de la variable $j$:
\[
	\sigma_j^2 = \frac{1}{n-1}\sum_{i=1}^n (x_{ij} - \mu_j)^2
\]

\paragraph{Covarianza} entre variables $j$ y $k$:
\[
	\sigma_{jk} = \frac{1}{n-1}\sum_{i=1}^n (x_{ij}-\mu_j)(x_{ik}-\mu_k)
\]

\paragraph{Matriz de covarianza}:
\[
	\mathbf{M}_x = \frac{1}{n-1}\mathbf{X}^\top\mathbf{X}
	\quad\in\R^{p\times p} \quad (\mathbf{X} \text{ centrada})
\]
$\mathbf{M}_x$ es simétrica semi-definida positiva; su diagonal contiene las varianzas
y sus fuera de la diagonal las covarianzas.

\subsubsection{Diagonalización y Autovectores}

La descomposición espectral de $\mathbf{M}_x$:
\[
	\mathbf{M}_x = \mathbf{P}\mathbf{D}\mathbf{P}^\top
\]
donde $\mathbf{P}$ es ortogonal (autovectores en columnas) y
$\mathbf{D} = \text{diag}(\lambda_1, \ldots, \lambda_p)$ con $\lambda_1 \geq \cdots \geq \lambda_p \geq 0$.

\begin{defbox}{Interpretación geométrica}
	\begin{itemize}
		\item Cada autovector $\mathbf{v}_j$ es una \textbf{dirección} en el espacio de features.
		\item El autovalor $\lambda_j$ es la \textbf{varianza} explicada en esa dirección.
		\item La primera CP ($\mathbf{v}_1$) apunta hacia la dirección de máxima varianza;
		      la segunda es perpendicular a la primera y tiene la segunda mayor varianza; etc.
	\end{itemize}
\end{defbox}

\begin{figure}[H]
	\centering
	\begin{tikzpicture}
		% Nube de puntos elíptica
		\draw[->] (-3,0) -- (3,0) node[right]{$x_1$};
		\draw[->] (0,-2) -- (0,2) node[above]{$x_2$};
		% Puntos
		\foreach \x/\y in {
				-2.0/-0.2, -1.8/0.1, -1.5/-0.3, -1.2/0.2,
				-0.8/-0.1, -0.3/0.3, 0/0, 0.4/-0.2,
				0.8/0.2, 1.2/0.1, 1.6/-0.1, 2.0/0.2}
		\fill[azul!60] ({0.7*\x-0.3*\y},{0.3*\x+0.7*\y}) circle (1.5pt);
		% PC1
		\draw[-Stealth, rojo, very thick] (0,0) --
		node[above right, font=\footnotesize]{PC1} (2.0,0.86);
		% PC2
		\draw[-Stealth, verde, very thick] (0,0) --
		node[above left, font=\footnotesize]{PC2} (-0.43,1.0);
	\end{tikzpicture}
	\caption{Las componentes principales son los ejes de la elipse de la nube de puntos.
		PC1 apunta a la dirección de mayor dispersión.}
\end{figure}

\subsubsection{Algoritmo PCA}

\begin{algbox}{PCA Paso a Paso}
	\begin{algorithmic}[1]
		\State Estandarizar cada feature: $x_{ij} \leftarrow (x_{ij} - \mu_j)/\sigma_j$
		\State Calcular $\mathbf{M}_x = \frac{1}{n-1}\mathbf{X}^\top\mathbf{X}$
		\State Descomponer: $\mathbf{M}_x \mathbf{v}_j = \lambda_j \mathbf{v}_j$
		\State Ordenar autovectores por autovalor descendente
		\State Seleccionar los primeros $k$ autovectores: $\mathbf{P}_k \in \R^{p\times k}$
		\State Proyectar datos: $\mathbf{Z} = \mathbf{X}\,\mathbf{P}_k \in \R^{n\times k}$
	\end{algorithmic}
\end{algbox}

\subsubsection{Varianza Explicada}

\[
	\text{Varianza explicada por CP}_j = \frac{\lambda_j}{\sum_{l=1}^p \lambda_l}
\]
\[
	\text{Varianza acumulada por primeras } k \text{ CPs} = \frac{\sum_{j=1}^k \lambda_j}{\sum_{l=1}^p \lambda_l}
\]

\begin{figure}[H]
	\centering
	\begin{tikzpicture}
		\begin{axis}[
				width=9cm, height=5cm,
				xlabel={Componente Principal},
				ylabel={\%},
				ymin=0, ymax=105,
				xtick={1,...,8},
				bar width=0.5cm,
				ymajorgrids, grid style={gray!20},
				title={Varianza Explicada (Scree Plot)},
				legend pos=outer north east
			]
			\addplot[ybar, fill=azul!50, draw=azul] coordinates {
					(1,40)(2,25)(3,15)(4,9)(5,5)(6,3)(7,2)(8,1)
				};
			\addplot[rojo, thick, mark=*] coordinates {
					(1,40)(2,65)(3,80)(4,89)(5,94)(6,97)(7,99)(8,100)
				};
			\addlegendentry{Var. exp. individual}
			\addlegendentry{Var. acumulada}
		\end{axis}
	\end{tikzpicture}
	\caption{Scree plot: se elige $k$ donde la curva acumulada supera el umbral deseado
		(usualmente 85–95\%).}
\end{figure}

\subsection{Métodos No Lineales de Reducción de Dimensionalidad}

PCA asume relaciones \textbf{lineales}; cuando los datos yacen en una variedad
no lineal, se necesitan métodos más flexibles.

\begin{figure}[H]
	\centering
	\begin{tikzpicture}
		\tikzset{box/.style={rectangle, rounded corners, draw=azul!60, fill=azul!8,
					text width=3cm, align=center, minimum height=1.2cm, font=\footnotesize}}
		\node[box] (pca) at (0,0)    {PCA\\{\tiny lineal, rápido, interpretable}};
		\node[box] (tsne) at (4.5,0) {t-SNE\\{\tiny no lineal, local, lento, sin proyección nueva}};
		\node[box] (umap) at (9,0)   {UMAP\\{\tiny no lineal, local+global, rápido, proyecta nuevas muestras}};
		\draw[Stealth-, gray!60, thick] (pca.east) -- (tsne.west)
		node[midway, above, font=\tiny]{complejidad};
		\draw[Stealth-, gray!60, thick] (tsne.east) -- (umap.west);
	\end{tikzpicture}
\end{figure}

\subsubsection{t-SNE}
\begin{itemize}
	\item Convierte distancias euclídeas en probabilidades de similitud (distribución Gaussiana en alta dimensión, t-Student en baja).
	\item Minimiza la divergencia KL entre distribuciones de alta y baja dimensión.
	\item Preserva la estructura local (vecinos cercanos); puede distorsionar la global.
	\item \textbf{No} es apto para proyectar nuevas muestras; sólo para visualización exploratoria.
\end{itemize}

\subsubsection{UMAP}
\begin{itemize}
	\item Basado en teoría de topología algebraica (variedades Riemannianas).
	\item Preserva estructura local \textbf{y} global mejor que t-SNE.
	\item Mucho más rápido que t-SNE; puede proyectar nuevas muestras.
	\item Útil tanto para visualización como preprocesamiento para ML.
\end{itemize}

\begin{center}
	\begin{tabular}{lcccc}
		\toprule
		                         & \textbf{PCA} & \textbf{t-SNE} & \textbf{UMAP} & \textbf{PaCMAP} \\
		\midrule
		Relaciones               & Lineal       & No lineal      & No lineal     & No lineal       \\
		Estructura local         & Media        & Alta           & Alta          & Alta            \\
		Estructura global        & Alta         & Baja           & Media         & Media-Alta      \\
		Velocidad                & Muy rápida   & Lenta          & Rápida        & Rápida          \\
		Proyecta nuevas muestras & Sí           & No             & Sí            & No              \\
		Interpretabilidad        & Alta         & Baja           & Baja          & Baja            \\
		\bottomrule
	\end{tabular}
\end{center}

% =============================================================
\section{Aprendizaje No Supervisado: Clustering}
% =============================================================

\begin{defbox}{Clustering}
	Dado un conjunto de muestras \emph{sin etiquetas}, encontrar grupos (\emph{clusters})
	tales que muestras dentro del mismo grupo sean similares entre sí y diferentes
	de muestras de otros grupos.
\end{defbox}

La reducción de dimensionalidad es frecuentemente un paso previo al clustering:
se proyecta a 2D/3D con PCA/UMAP y se visualizan agrupamientos emergentes.

% =============================================================
\section{Resumen Comparativo de Modelos}
% =============================================================

\begin{center}
	\begin{tabular}{lcccc}
		\toprule
		\textbf{Modelo}              & \textbf{Sesgo} & \textbf{Varianza} & \textbf{Interpretable} & \textbf{No lineal} \\
		\midrule
		Regresión Lineal/Logística   & Alto           & Bajo              & Muy alta               & No                 \\
		Árbol de decisión (profundo) & Bajo           & Alto              & Media                  & Sí                 \\
		Árbol de decisión (podado)   & Medio          & Medio             & Alta                   & Sí                 \\
		Bagging / RF                 & Bajo           & Bajo              & Baja                   & Sí                 \\
		Gradient Boosting            & Muy bajo       & Bajo              & Baja                   & Sí                 \\
		\bottomrule
	\end{tabular}
\end{center}

\subsection{Bias-Variance Trade-off}

\[
	\E\!\left[(y - \hat{f}(x))^2\right]
	= \underbrace{\left(\E[\hat{f}(x)] - f(x)\right)^2}_{\text{Sesgo}^2}
	+ \underbrace{\E\!\left[\left(\hat{f}(x) - \E[\hat{f}(x)]\right)^2\right]}_{\text{Varianza}}
	+ \underbrace{\sigma^2_\varepsilon}_{\text{Ruido irreducible}}
\]

\begin{figure}[H]
	\centering
	\begin{tikzpicture}
		\begin{axis}[
				width=9cm, height=5.5cm,
				xlabel={Complejidad del modelo},
				ylabel={Error},
				xmin=0, xmax=10, ymin=0, ymax=2,
				grid=both, grid style={gray!15},
				title={Bias-Variance Trade-off},
				legend pos=north east
			]
			\addplot[rojo, thick, domain=0.5:10, samples=80]
			{3.5/(x+0.5) + 0.1};
			\addlegendentry{Sesgo²}
			\addplot[azul, thick, domain=0.5:10, samples=80]
			{0.005*x^2.2 + 0.05};
			\addlegendentry{Varianza}
			\addplot[verde, very thick, domain=0.5:10, samples=80]
			{3.5/(x+0.5) + 0.005*x^2.2 + 0.15};
			\addlegendentry{Error total}
			\draw[dashed, naranja, thick] (axis cs:5.5,0) -- (axis cs:5.5,2);
			\node[naranja, font=\tiny] at (axis cs:7,1.8) {Complejidad óptima};
		\end{axis}
	\end{tikzpicture}
	\caption{El error total tiene un mínimo en la complejidad óptima del modelo.}
\end{figure}

\subsection{Cuándo Usar Cada Modelo}

\begin{center}
	\begin{tabular}{p{3.5cm}p{9cm}}
		\toprule
		\textbf{Modelo}         & \textbf{Cuándo usarlo}                          \\
		\midrule
		Reg. Lineal / Logística &
		Cuando se requiere interpretabilidad, las relaciones son lineales
		o los datos son escasos.                                                  \\
		Árbol de decisión       &
		Para exploración rápida, cuando la interpretabilidad es crítica, o
		como base para ensambles.                                                 \\
		Random Forest           &
		Buen punto de partida robusto; menos sensible a hiperparámetros que GB.   \\
		Gradient Boosting       &
		Cuando la performance predictiva es prioritaria; tabulares estructurados. \\
		PCA                     &
		Preprocesamiento, visualización, eliminar correlaciones, ruido.           \\
		UMAP / t-SNE            &
		Visualización exploratoria de datos de alta dimensión.                    \\
		\bottomrule
	\end{tabular}
\end{center}

% =============================================================
\appendix
\section{Formulario Rápido}
% =============================================================

\begin{multicols}{2}
	\subsection*{Regresión Lineal}
	\[
		\hat{\bm{\beta}} = (\mathbf{X}^\top\mathbf{X})^{-1}\mathbf{X}^\top\mathbf{y}
	\]

	\subsection*{Logistic}
	\[
		\hat{p} = \sigma(\bm{\beta}^\top\mathbf{x}) = \frac{1}{1+e^{-\bm{\beta}^\top\mathbf{x}}}
	\]

	\subsection*{Lasso}
	\[
		\min_{\mathbf{w}}\; J(\mathbf{w}) + \lambda\|\mathbf{w}\|_1
	\]

	\subsection*{Ridge}
	\[
		\min_{\mathbf{w}}\; J(\mathbf{w}) + \lambda\|\mathbf{w}\|_2^2
	\]

	\subsection*{Gini}
	\[
		\text{Gini}_m = 1 - \sum_k p_{mk}^2
	\]

	\subsection*{Entropía}
	\[
		H_m = -\sum_k p_{mk}\log_2 p_{mk}
	\]

	\subsection*{AdaBoost peso modelo}
	\[
		\alpha_m = \frac{1}{2}\ln\frac{1-\varepsilon_m}{\varepsilon_m}
	\]

	\subsection*{Gradient Boosting}
	\[
		r_{im} = -\frac{\partial L(y_i, F_{m-1}(x_i))}{\partial F_{m-1}(x_i)}
	\]

	\subsection*{PCA proyección}
	\[
		\mathbf{Z} = \mathbf{X}\,\mathbf{P}_k
	\]

	\subsection*{Varianza PCA}
	\[
		\text{VE}_j = \frac{\lambda_j}{\sum_l \lambda_l}
	\]

	\subsection*{Bias-Variance}
	\[
		\E[(y-\hat{f})^2] = \text{Sesgo}^2 + \text{Var} + \sigma^2
	\]

	\subsection*{OOB estimate}
	Cada muestra es OOB para $\approx e^{-1} \approx 36.8\%$ de árboles.
\end{multicols}

\end{document}
